% !PRINTED 609
\weberpage{609}
In demselben Falle finden sich die cyklischen und überhaupt alle irreduciblen
\gesperrt{Abel}'schen Gleichungen von ungeradem Grade, die also niemals durch reelle
Radicale lösbar sind.

Hiernach können wir z. B. die cubischen Gleichungen im Körper der rationalen Zahlen in
drei oder vier Arten unterscheiden, die alle in gewissen von Alters her berühmten
geometrischen Problemen auftreten. Wir sehen dabei von den reduciblen Gleichungen ab.
Der Grad der \gesperrt{Galois}'schen Gruppe muss dann immer durch 3 theilbar sein
(\S.\,154, 7.) und ist also entweder gleich 3 oder gleich 6. Die Gruppe kann also nach
\S.\,157 niemals durch Adjunction von blossen Quadratwurzeln auf einen niedrigeren als
den dritten Grad reducirt werden; also kann auch die Gleichung nicht durch
Quadratwurzeln gelöst werden. Die entsprechenden geometrischen Probleme sind nicht mit
Cirkel und Lineal zu lösen.

Ist die Gruppe vom Grade 3, so haben wir eine cyklische Gleichung. Hierher gehören die
aus der Kreistheilung stammenden Gleichungen, z. B. die, von der die Construction des
regelmässigen Siebenecks abhängt. Die drei Wurzeln einer solchen Gleichung sind reell
und können nicht durch reelle Radicale ausgedrückt werden. Unter den cubischen
Gleichungen mit einer Gruppe $6^{\text{ten}}$ Grades sind zu unterscheiden die mit
positiver und mit negativer Discriminante.

Die ersten gehören zum casus irreducibilis und lassen sich zurückführen auf die
Gleichung, von der die Dreitheilung eines beliebigen Winkels abhängt. Zu den
Gleichungen mit negativer Discriminante gehören als specielle Fälle auch die reinen
cubischen Gleichungen $x^3 = a$, wenn $a$ keine Cubikzahl ist. Für $a = 2$ ergiebt sich
die Gleichung, von der das \gesperrt{Delische Problem} der Würfelverdoppelung abhängt.

{\centering\S.\,180.\par}
{\centering \gesperrt{Metacyklische Gleichungen von Primzahlgrad.}\par}

Die allgemeinen Bedingungen, die wir im \S.\,176 gefunden haben, sind noch nicht
einfach genug, um eine unmittelbare Anwendung auf die Ermittelung von metacyklischen
Gleichungen oder auf die Entscheidung über die Lösbarkeit einer vorgelegten
%CONT>>
